% GENERATED FILE. Edit canonical lessons, then run npm run generate:books.
% canonical-source-hash: fnv1a64:6b2ebea1a1a28d8c
% canonical-lessons: BN-C328-prishtha, BN-C328-onubad-kora, BN-C328-pandulipi, BN-C328-sarangsho, BN-C328-rochona-kora

\chapter{Inside a Book}
\label{ch:bn-inside-a-book}
\hlchaptermodality{328}

\begin{chapteropening}
\textbf{By the end of this chapter:} \emph{I can talk about pages, translations, manuscripts and summaries.}

Complete the last lesson of chapter 328: I can talk about pages, translations, manuscripts and summaries.
\end{chapteropening}

\section[\texorpdfstring{\bn{পৃষ্ঠা} (prishṭhā)}{পৃষ্ঠা (prishṭhā)}]{\bn{পৃষ্ঠা} (prishṭhā) — a page}

\label{lesson:BN-C328-prishtha}



\begin{quote}
Before the new one: say the Bengali for a chapter of a book, then the Bengali for a publisher.
\end{quote}

\subsection*{You'll want to know: \bn{পৃষ্ঠা}}
\textbf{\bn{পৃষ্ঠা}} — \emph{prishṭhā} — \textquotedblleft{}a page\textquotedblright{}.

\begin{quote}
\bn{বইটার} \bn{প্রথম} \bn{পৃষ্ঠায়} \bn{লেখকের} \bn{ছবি} \bn{আছে।}
\end{quote}

\begin{quote}
\emph{bôiṭār prothom prishṭhāy lekhôker chhôbi āchhe.}
\end{quote}

\begin{quote}
There is a picture of the writer on the first page of the book.
\end{quote}

\begin{cousinweb}
\textbf{\bn{পৃষ্ঠা}} (\emph{prishṭhā}) is Sanskrit \emph{pṛṣṭha}, 'the back; the upper side, a surface' (Monier-Williams); the sense 'a page', one face of a leaf, is the later one. Everyday speech also says \textbf{\bn{পাতা}}.
\end{cousinweb}

\begin{grammarlens}[title={What you've built}]
One more word for telling what happened.
\end{grammarlens}

\subsection*{Your turn}
\begin{itemize}
\raggedright
  \item \emph{Say it:} \emph{prishṭhā}
  \item \emph{Say it:} \emph{prishṭhā}, once more
  \item \emph{Say it:} \emph{bôiṭār prothom prishṭhāy lekhôker chhôbi āchhe.} — There is a picture of the writer on the first page of the book.
  \item \emph{Recall:} say the Bengali for a chapter of a book, then the Bengali for a publisher, then say \emph{prishṭhā} again
\end{itemize}

\begin{tcolorbox}[breakable,colback=teal!4,colframe=teal!35!black,title={Before you move on}]
What does \bn{পৃষ্ঠা} mean? (\textquotedblleft{}A page\textquotedblright{}.) Say it once more.
\end{tcolorbox}
\section[\texorpdfstring{\bn{অনুবাদ} \bn{করা} (ônubād kôrā)}{অনুবাদ করা (ônubād kôrā)}]{\bn{অনুবাদ} \bn{করা} (ônubād kôrā) — to translate}

\label{lesson:BN-C328-onubad-kora}



\begin{quote}
Before the new one: say the Bengali for a publisher, then the Bengali for a page.
\end{quote}

\subsection*{You'll want to know: \bn{অনুবাদ} \bn{করা}}
\textbf{\bn{অনুবাদ} \bn{করা}} — \emph{ônubād kôrā} — \textquotedblleft{}to translate\textquotedblright{}.

\begin{quote}
\bn{তিনি} \bn{বাংলা} \bn{উপন্যাসটা} \bn{ইংরেজিতে} \bn{অনুবাদ} \bn{করলেন।}
\end{quote}

\begin{quote}
\emph{tini bānglā upônyāsṭā ingrejite ônubād kôrlen.}
\end{quote}

\begin{quote}
He translated the Bengali novel into English.
\end{quote}

Past (he/she …): \bn{তিনি} \bn{কবিতাটা} \bn{অনুবাদ} \bn{করলেন}, \emph{tini kôbitāṭā ônubād kôrlen}.

\begin{cousinweb}
\textbf{\bn{অনুবাদ} \bn{করা}} (\emph{ônubād kôrā}) puts Sanskrit \emph{anuvāda}, 'saying after, repeating; translation' (Monier-Williams), before \bn{করা}: to say again in another language.
\end{cousinweb}

\begin{grammarlens}[title={What you've built}]
One more word for telling what happened.
\end{grammarlens}

\subsection*{Your turn}
\begin{itemize}
\raggedright
  \item \emph{Say it:} \emph{ônubād kôrā}
  \item \emph{Say it:} \emph{ônubād kôrā}, once more
  \item \emph{Say it:} \emph{tini bānglā upônyāsṭā ingrejite ônubād kôrlen.} — He translated the Bengali novel into English.
  \item \emph{Recall:} say the Bengali for a publisher, then the Bengali for a page, then say \emph{ônubād kôrā} again
\end{itemize}

\begin{tcolorbox}[breakable,colback=teal!4,colframe=teal!35!black,title={Before you move on}]
What does \bn{অনুবাদ} \bn{করা} mean? (\textquotedblleft{}To translate\textquotedblright{}.) Say it once more.
\end{tcolorbox}
\section[\texorpdfstring{\bn{পাণ্ডুলিপি} (pānḍulipi)}{পাণ্ডুলিপি (pānḍulipi)}]{\bn{পাণ্ডুলিপি} (pānḍulipi) — a manuscript}

\label{lesson:BN-C328-pandulipi}



\begin{quote}
Before the new one: say the Bengali for a page, then the Bengali for to translate.
\end{quote}

\subsection*{You'll want to know: \bn{পাণ্ডুলিপি}}
\textbf{\bn{পাণ্ডুলিপি}} — \emph{pānḍulipi} — \textquotedblleft{}a manuscript\textquotedblright{}.

\begin{quote}
\bn{নাতি} \bn{দাদুর} \bn{বাক্সে} \bn{কবিতার} \bn{পাণ্ডুলিপি} \bn{পেল।}
\end{quote}

\begin{quote}
\emph{nāti dādur bākse kôbitār pānḍulipi pelô.}
\end{quote}

\begin{quote}
The grandson found the manuscript of the poems in his grandfather's box.
\end{quote}

\begin{cousinweb}
\textbf{\bn{পাণ্ডুলিপি}} (\emph{pānḍulipi}) joins Sanskrit \emph{pāṇḍu}, 'pale, yellowish white' (Monier-Williams), to \emph{lipi}, 'writing'. Monier-Williams has the older \emph{pāṇḍulekha}, a first sketch in chalk; \bn{পাণ্ডুলিপি} is a work in the writer's own hand, before it is printed.
\end{cousinweb}

\begin{grammarlens}[title={What you've built}]
One more word for telling what happened.
\end{grammarlens}

\subsection*{Your turn}
\begin{itemize}
\raggedright
  \item \emph{Say it:} \emph{pānḍulipi}
  \item \emph{Say it:} \emph{pānḍulipi}, once more
  \item \emph{Say it:} \emph{nāti dādur bākse kôbitār pānḍulipi pelô.} — The grandson found the manuscript of the poems in his grandfather's box.
  \item \emph{Recall:} say the Bengali for a page, then the Bengali for to translate, then say \emph{pānḍulipi} again
\end{itemize}

\begin{tcolorbox}[breakable,colback=teal!4,colframe=teal!35!black,title={Before you move on}]
What does \bn{পাণ্ডুলিপি} mean? (\textquotedblleft{}A manuscript\textquotedblright{}.) Say it once more.
\end{tcolorbox}
\section[\texorpdfstring{\bn{সারাংশ} (sārāngshô)}{সারাংশ (sārāngshô)}]{\bn{সারাংশ} (sārāngshô) — a summary}

\label{lesson:BN-C328-sarangsho}



\begin{quote}
Before the new one: say the Bengali for to translate, then the Bengali for a manuscript.
\end{quote}

\subsection*{You'll want to know: \bn{সারাংশ}}
\textbf{\bn{সারাংশ}} — \emph{sārāngshô} — \textquotedblleft{}a summary\textquotedblright{}.

\begin{quote}
\bn{ক্লাসে} \bn{আমরা} \bn{গল্পটার} \bn{সারাংশ} \bn{লিখলাম।}
\end{quote}

\begin{quote}
\emph{klāse āmrā gôlpoṭār sārāngshô likhlām.}
\end{quote}

\begin{quote}
In class we wrote a summary of the story.
\end{quote}

\begin{cousinweb}
\textbf{\bn{সারাংশ}} (\emph{sārāngshô}) joins Sanskrit \emph{sāra}, 'the core, the essence', to \emph{aṃśa}, 'a share, a part' (Monier-Williams): the part that holds the core.
\end{cousinweb}

\begin{grammarlens}[title={What you've built}]
One more word for telling what happened.
\end{grammarlens}

\subsection*{Your turn}
\begin{itemize}
\raggedright
  \item \emph{Say it:} \emph{sārāngshô}
  \item \emph{Say it:} \emph{sārāngshô}, once more
  \item \emph{Say it:} \emph{klāse āmrā gôlpoṭār sārāngshô likhlām.} — In class we wrote a summary of the story.
  \item \emph{Recall:} say the Bengali for to translate, then the Bengali for a manuscript, then say \emph{sārāngshô} again
\end{itemize}

\begin{tcolorbox}[breakable,colback=teal!4,colframe=teal!35!black,title={Before you move on}]
What does \bn{সারাংশ} mean? (\textquotedblleft{}A summary\textquotedblright{}.) Say it once more.
\end{tcolorbox}
\section[\texorpdfstring{\bn{রচনা} \bn{করা} (rôchônā kôrā)}{রচনা করা (rôchônā kôrā)}]{\bn{রচনা} \bn{করা} (rôchônā kôrā) — to compose; to author}

\label{lesson:BN-C328-rochona-kora}



\begin{quote}
Before the new one: say the Bengali for a manuscript, then the Bengali for a summary.
\end{quote}

\subsection*{You'll want to know: \bn{রচনা} \bn{করা}}
\textbf{\bn{রচনা} \bn{করা}} — \emph{rôchônā kôrā} — \textquotedblleft{}to compose; to author\textquotedblright{}.

\begin{quote}
\bn{কবি} \bn{এই} \bn{গানটা} \bn{একটা} \bn{নৌকায়} \bn{বসে} \bn{রচনা} \bn{করলেন।}
\end{quote}

\begin{quote}
\emph{kôbi ei gānṭā ekṭā noukāy bôse rôchônā kôrlen.}
\end{quote}

\begin{quote}
The poet composed this song sitting in a boat.
\end{quote}

Past (he/she …): \bn{তিনি} \bn{একটা} \bn{কবিতা} \bn{রচনা} \bn{করলেন}, \emph{tini ekṭā kôbitā rôchônā kôrlen}.

\begin{cousinweb}
\textbf{\bn{রচনা} \bn{করা}} (\emph{rôchônā kôrā}) puts Sanskrit \emph{racanā}, 'arrangement, preparation; a literary work' (Monier-Williams), before \bn{করা}: to arrange words into a poem or a song.
\end{cousinweb}

\begin{grammarlens}[title={What you've built}]
One more word for telling what happened.
\end{grammarlens}

\subsection*{Your turn}
\begin{itemize}
\raggedright
  \item \emph{Say it:} \emph{rôchônā kôrā}
  \item \emph{Say it:} \emph{rôchônā kôrā}, once more
  \item \emph{Say it:} \emph{kôbi ei gānṭā ekṭā noukāy bôse rôchônā kôrlen.} — The poet composed this song sitting in a boat.
  \item \emph{Recall:} say the Bengali for a manuscript, then the Bengali for a summary, then say \emph{rôchônā kôrā} again
\end{itemize}

\begin{tcolorbox}[breakable,colback=teal!4,colframe=teal!35!black,title={Before you move on}]
What does \bn{রচনা} \bn{করা} mean? (\textquotedblleft{}To compose; to author\textquotedblright{}.) Say it once more.
\end{tcolorbox}
